% =====================================================================
% §6 Discussion + Conclusion + Climate impact
% =====================================================================

\section{Discussion}
\label{sec:discussion}

\paragraph{Implications for retraining cadence.}
A three-year retraining gap visibly hurts in this dataset. Per-zone
MAPE in May~2026 ranges from 7~\% (NEMA-Boston) to 78~\% (WCMA) for the
trained baseline---the headline 2022 number masks zone-specific
catastrophic behavior. Reporting only the overall MAPE on a 2-day test
slice (5.24~\%) gave us essentially no warning signal. Per-zone
surveillance plus a continuously-refreshed validation slice (the
deployed system has both) catches drift earlier.

\paragraph{Per-zone $\alpha_z$ as drift localiser.}
The zones the validation-tuned $\alpha_z$ flagged in 2022 as
``low-baseline-trust'' ($\alpha = 0$ in CT, SEMA, NEMA-Boston) overlap
substantially with the zones suffering largest 2026 drift. With $n =
8$, this is suggestive rather than conclusive, but it points at a
practical use case: weighted-fusion weights tuned even before
deployment may carry per-zone-specific information about where the
model is fragile.

\paragraph{Limitations.}
(i) Two windows are better than one but not as strong as a full year
of monthly snapshots; we plan to extend the rolling backtest to a
calendar of windows in future work.
(ii) The BTM-solar correlation is observational at $n = 5$ states. We
cannot claim causation without a controlled counterfactual study.
(iii) The deployed model uses HRRR forecasts (not analyses) for the
future window; this is a known training-time vs deployment-time
distribution mismatch quantified in Appendix~\ref{app:forecast_shift}
(small effect, but non-zero).
(iv) Single ISO. Replicating this experiment under CAISO, ERCOT, or
PJM would test whether the mechanism generalises, but each ISO has its
own data-access patterns and would require building a similar fetch
stack.

\paragraph{Climate impact and broader use.}
Day-ahead demand forecasting is a load-bearing input to grid
operations. Our deployed model is far below production-grade and should
not be used for real dispatch. The intended audience for the
deployment-drift case-study finding is researchers and grid forecasters
who own production models trained pre-2022 and need a calibrated
expectation of how BTM solar buildout changes their per-zone error
distribution. Demand-response programs and renewable-energy integration
analyses use forecasts of this kind; better-calibrated drift estimates
help integrate distributed solar more reliably.

\paragraph{Conclusion.}
A multimodal CNN-Transformer + Chronos-Bolt-mini ensemble for ISO-NE
day-ahead forecasting reaches 4.21~\% MAPE on a 2-day 2022 test slice.
The same trained pipeline, deployed publicly and re-evaluated on
rolling 7-day windows in May~2025 and May~2026, drifts in a
zone-specific pattern that aligns with state-level BTM solar density.
The pipeline-equivalence experiment, the bootstrap-CI'd headline
numbers, and the daily-refreshed live demo make the case study
reproducible end-to-end. The artifacts are public
(Appendix~\ref{app:repro}).
